% PDF title page (latex_elements["maketitle"] in conf.py)
\begin{titlepage}
\begin{tikzpicture}[remember picture,overlay]
  % black band with orange rule
  \fill[novanero] (current page.north west) rectangle ([yshift=-14cm]current page.north east);
  \fill[novaarancio] ([yshift=-14cm]current page.north west) rectangle ([yshift=-14.25cm]current page.north east);
  % motif: the three ring segments of the NOVA symbol
  \begin{scope}[shift={([xshift=-4.6cm,yshift=-7.2cm]current page.north east)}]
    \foreach \a in {90,210,330}{
      \draw[novaarancio,line width=9pt,line cap=round]
        ([shift={(\a-38:2.3cm)}]0,0) arc[start angle=\a-38,end angle=\a+38,radius=2.3cm];
      \draw[white,opacity=0.35,line width=0.7pt]
        ([shift={(\a+60:1.35cm)}]0,0) arc[start angle=\a+60,end angle=\a+120,radius=1.35cm];
    }
  \end{scope}
\end{tikzpicture}
\vspace*{0.6cm}
{\novasans\color{white!65!novanero}\small\bfseries\addfontfeatures{LetterSpace=18}TECHNICAL MANUAL\par}
\vspace{10pt}
\includegraphics[height=1.75cm]{marchio-bianco.pdf}\par
\vspace{12pt}
{\novasans\color{white}\LARGE NOVA documentation\par}
\vspace{6pt}
{\novasans\color{white!65!novanero}\large How it works\par}
\vspace{8.4cm}
\begin{minipage}{0.84\textwidth}
\large\color{novainchiostro}
This document describes how NOVA works: architecture, rotor models, block workflows, and the
verification and validation process.
\textbf{[TO COMPLETE]} Replace this paragraph with an introduction to the manual: who it is for,
what it covers, and which software version it refers to.
\end{minipage}
\vfill
{\novasans\color{novagrigio}\small NOVA \hfill \textbf{[TO COMPLETE]} version and date}
\end{titlepage}
